\documentclass[11pt,a4paper]{article}

\usepackage[utf8]{inputenc}
\usepackage[T1]{fontenc}
\usepackage[french]{babel}

\usepackage{styles/couleurs}
\usepackage{styles/mise_en_page}
\usepackage{styles/boites}
\usepackage{styles/maths}
\usepackage{styles/figures}

\usepackage{amsmath}
\usepackage{amssymb}
\usepackage[most]{tcolorbox}

\geometry{
  top=10mm,
  bottom=10mm,
  left=12mm,
  right=12mm
}

\pagestyle{empty}
\setlength{\parindent}{0pt}
\setlength{\parskip}{0pt}

% =========================================================
% BOÎTE EXERCICE
% =========================================================

\newtcolorbox{exercicebox}[1]{
  enhanced,
  colback=white,
  colframe=ActionC!45,
  boxrule=0.6pt,
  arc=1.5mm,
  left=4mm,
  right=4mm,
  top=2mm,
  bottom=2mm,
  title={#1},
  fonttitle=\sffamily\bfseries,
  coltitle=ActionC,
  colbacktitle=white,
  boxed title style={
    colframe=white,
    colback=white
  },
  attach boxed title to top left={
    xshift=2mm,
    yshift=-1mm
  }
}

\begin{document}

% =========================================================
% TITRE
% =========================================================

\begin{center}
{\sffamily\LARGE\bfseries\color{TitreC}
AP -- Je rédige avec le théorème de Pythagore}
\end{center}

\vspace{1mm}

\begin{tcolorbox}[
  enhanced,
  colback=SavoirBg,
  colframe=SavoirC!50,
  boxrule=0.6pt,
  arc=1.5mm,
  left=4mm,
  right=4mm,
  top=2mm,
  bottom=2mm
]

Pour chaque situation :

\medskip

\textbf{1.} Fais un schéma à main levée et code l'angle droit.

\textbf{2.} Rédige toutes les étapes du théorème de Pythagore.

\end{tcolorbox}

\vspace{3mm}

% =========================================================
% LIGNE 1
% =========================================================

\noindent
\begin{minipage}[t]{0.485\linewidth}

\begin{exercicebox}{Triangle ABC}
Le triangle \(ABC\) est rectangle en \(A\).

\[
AB=6\text{ cm}
\qquad
AC=9\text{ cm}
\]

Calculer \(BC\).

\textbf{Donner une valeur approchée au dixième.}
\end{exercicebox}

\end{minipage}
\hfill
\begin{minipage}[t]{0.485\linewidth}

\begin{exercicebox}{Triangle DEF}
Le triangle \(DEF\) est rectangle en \(E\).

\[
DF=14\text{ cm}
\qquad
DE=5\text{ cm}
\]

Calculer \(EF\).

\textbf{Donner une valeur approchée au centième.}
\end{exercicebox}

\end{minipage}

\vspace{3mm}

% =========================================================
% LIGNE 2
% =========================================================

\noindent
\begin{minipage}[t]{0.485\linewidth}

\begin{exercicebox}{Triangle GHI}
Le triangle \(GHI\) est rectangle en \(H\).

\[
GH=7\text{ cm}
\qquad
HI=12\text{ cm}
\]

Calculer \(GI\).

\textbf{Donner une valeur approchée au centième.}
\end{exercicebox}

\end{minipage}
\hfill
\begin{minipage}[t]{0.485\linewidth}

\begin{exercicebox}{Triangle JKL}
Le triangle \(JKL\) est rectangle en \(K\).

\[
JL=18\text{ cm}
\qquad
JK=7\text{ cm}
\]

Calculer \(KL\).

\textbf{Donner une valeur approchée au dixième.}
\end{exercicebox}

\end{minipage}

\vspace{3mm}

% =========================================================
% LIGNE 3
% =========================================================

\noindent
\begin{minipage}[t]{0.485\linewidth}

\begin{exercicebox}{Triangle MNP}
Le triangle \(MNP\) est rectangle en \(N\).

\[
MN=8\text{ cm}
\qquad
NP=13\text{ cm}
\]

Calculer \(MP\).

\textbf{Donner une valeur approchée au dixième.}
\end{exercicebox}

\end{minipage}
\hfill
\begin{minipage}[t]{0.485\linewidth}

\begin{exercicebox}{Triangle QRS}
Le triangle \(QRS\) est rectangle en \(Q\).

\[
RS=19\text{ cm}
\qquad
QS=12\text{ cm}
\]

Calculer \(QR\).

\textbf{Donner une valeur approchée au centième.}
\end{exercicebox}

\end{minipage}

\vspace{3mm}

% =========================================================
% LIGNE 4
% =========================================================

\noindent
\begin{minipage}[t]{0.485\linewidth}

\begin{exercicebox}{Triangle TUV}
Le triangle \(TUV\) est rectangle en \(V\).

\[
TV=9\text{ cm}
\qquad
UV=14\text{ cm}
\]

Calculer \(TU\).

\textbf{Donner une valeur approchée au centième.}
\end{exercicebox}

\end{minipage}
\hfill
\begin{minipage}[t]{0.485\linewidth}

\begin{exercicebox}{Triangle XYZ}
Le triangle \(XYZ\) est rectangle en \(Y\).

\[
XZ=20\text{ cm}
\qquad
XY=7\text{ cm}
\]

Calculer \(YZ\).

\textbf{Donner une valeur approchée au dixième.}
\end{exercicebox}

\end{minipage}

\end{document}